\section{Conclusion and Outlook}

This work has examined the question of photon mass from complementary 
perspectives: gauge symmetry, field theory, laboratory experiments, 
astrophysical observations, and the helicity structure of quantum 
entanglement. Together, these independent lines of reasoning lead to a 
coherent and remarkably robust conclusion: if the photon has a mass at 
all, it must be extraordinarily small—far below any currently 
detectable threshold.

From a theoretical standpoint, the masslessness of the photon is deeply 
rooted in the structure of electromagnetism. Gauge symmetry forbids a 
mass term, ensures charge conservation, and guarantees the 
renormalizability of quantum electrodynamics. A non-zero photon mass 
would require abandoning the unbroken $U(1)$ symmetry and introducing 
new physics beyond the Standard Model.

From an empirical standpoint, the absence of deviations from Maxwell’s 
theory across scales ranging from laboratory experiments to galactic 
magnetic fields constrains the photon mass to values as small as
\[
m_\gamma \lesssim 10^{-27}\,\text{eV}.
\]
Astrophysical observations set the strongest limits, while laboratory 
tests provide independent confirmation under controlled and 
model-independent conditions.

An additional and conceptually powerful indicator comes from quantum 
entanglement: the experimentally observed two-state helicity structure 
of photons is consistent only with a massless spin-1 particle. A massive 
photon would possess a third, longitudinal degree of freedom, which 
would necessarily manifest itself in polarization correlations and 
entanglement experiments. No such deviations have ever been observed.

Taken together, these arguments strongly support the view that the 
photon is massless in a fundamental sense. The distinction between 
“almost zero” and “exactly zero” continues to motivate both conceptual 
reflection and experimental innovation. While experiments can never 
prove that a quantity is exactly zero, the convergence of theoretical 
consistency and empirical limits overwhelmingly favours exact 
masslessness.

\subsection*{Outlook}

Future observational and theoretical developments will continue to 
refine our understanding:

\begin{itemize}
	\item Improved time-of-flight measurements from high-energy 
	astrophysical sources may further tighten bounds on dispersion effects.
	\item Precision magnetometry in future space missions could probe 
	deviations from Maxwellian behaviour at solar-system scales.
	\item Advances in quantum technologies may enable new tests of photon 
	helicity and entanglement over long distances.
	\item Theoretical progress in quantum gravity and beyond-Standard-Model 
	frameworks may offer deeper insight into why certain particles—
	including the photon—are fundamentally massless.
\end{itemize}

The question of photon mass thus remains more than a numerical 
determination; it serves as a window into the structure of physical law.  
The current evidence overwhelmingly supports the masslessness of the 
photon, yet the conceptual richness of the problem continues to inspire 
both theoretical investigation and experimental exploration.
